% Prostudováno: ucebnice/02_Funkce_a_rovnice/Linearni_funkce_a_rovnice/Linerání rovnice.tex,
% Soustavy lineárních rovnic, geometrický význam a počet řešení (ř. 998–1122).
% Realisticky.cz: externi_zdroje/realisticky/02_Funkce_a_rovnice/03_Nektere_rovnice_a_nerovnice_prevoditelne_na_linearni/020320_Graficke_reseni_soustav_linearnich_rovnic_a_nerovnic/020320_Lekce.pdf
% a 020311_Soustavy_dvou_linearnich_rovnic_o_dvou_neznamych_II/020311_Lekce.pdf; využito propojení algebraického výsledku 0=0/0=c s polohou přímek. Grafy a úlohy vlastní.
\documentclass[11pt,a4paper]{article}
% Spolecny styl pro ucitelske pripravy MAT (Work/MAT)
% Verze 1.1, 2026-10-23
\NeedsTeXFormat{LaTeX2e}
\ProvidesPackage{priprava}[2026/10/23 v1.1 Sablona ucitelskych priprav MAT]

% Zakladni typografie a cestina
\RequirePackage[T1]{fontenc}
\RequirePackage[utf8]{inputenc}
\RequirePackage[czech]{babel}
\RequirePackage{lmodern}
\RequirePackage{microtype}
\RequirePackage{amsmath,amssymb}
\RequirePackage{xcolor}
\RequirePackage{enumitem}
\RequirePackage[most]{tcolorbox}
\RequirePackage{tabularx,booktabs,array}
\RequirePackage{fancyhdr}
\RequirePackage{needspace}
\RequirePackage{geometry}
\RequirePackage[hidelinks]{hyperref}

\geometry{a4paper,top=17mm,bottom=18mm,left=23mm,right=23mm}
\definecolor{notegray}{gray}{0.39}
\setlength{\parindent}{0pt}
\setlength{\parskip}{2.5pt plus 1pt minus 0.5pt}
\setlist[itemize]{topsep=2pt,itemsep=1pt,parsep=0pt,leftmargin=1.45em}
\setlist[enumerate]{topsep=2pt,itemsep=2pt,parsep=0pt,leftmargin=1.65em}
\setcounter{secnumdepth}{0}
\renewcommand{\arraystretch}{1.12}

% Zahlavi: vlevo rocnik | globalni tema, vpravo cislo: nazev hodiny
% Zapati zustava zcela prazdne.
\pagestyle{fancy}
\fancyhf{}
\renewcommand{\headrulewidth}{0.35pt}
\renewcommand{\footrulewidth}{0pt}
\setlength{\headheight}{14pt}
\newcommand{\nastavPripravu}[4]{%
  \fancyhead[L]{\small #1\enspace\textbar\enspace #2}%
  \fancyhead[R]{\small #3: #4}%
}

% Minimalisticke cernobile boxy. Kazda definice zvlast.
\tcbset{commonlesson/.style={enhanced,breakable=false,colback=white,colframe=black,
  colbacktitle=white,coltitle=black,boxrule=0.45pt,
  arc=0.8mm,outer arc=0.8mm,left=3mm,right=3mm,top=1.6mm,bottom=1.6mm,
  toptitle=1mm,bottomtitle=0.7mm,fonttitle=\bfseries,
  attach boxed title to top left={xshift=2.5mm,yshift=-1.15mm},
  boxed title style={colback=white,colframe=white,boxrule=0pt,left=1mm,right=1mm}}}
\newtcolorbox{zadaniBox}[1][]{commonlesson,title={Zadání},#1}
\newtcolorbox{prikladBox}[1][]{commonlesson,title={Příklad},#1}
\newtcolorbox{definiceBox}[1][]{commonlesson,title={Definice},#1}

% Spolecny box pro zadani a jeho reseni, bez mezery mezi castmi.
% Pouziti:
%   \begin{prikladReseniBox}
%     Zadani ulohy.
%     \tcblower
%     \begin{align*} ... \end{align*}
%   \end{prikladReseniBox}
% Spodni cast ma jemny sedy podklad a delici linku; lze v ni sazet
% vice radku matematiky. Obe casti maji jen jeden spolecny ramecek.
\tcbset{solutionlesson/.style={
  commonlesson,
  enhanced jigsaw,
  breakable,
  colbacklower=black!3,
  segmentation style={solid,draw=black!25,line width=0.35pt},
  fontlower=\small\color{notegray},
  before lower={\noindent\textbf{Řešení:}\enspace},
  leftlower=3mm,rightlower=3mm
}}
\newtcolorbox{prikladReseniBox}[1][]{solutionlesson,title={Příklad},#1}
\newtcolorbox{zadaniReseniBox}[1][]{solutionlesson,title={Zadání},#1}

% Pomucky pro autora pripravy: sedive poznamky; \reseni pro starsi pripravy
\newcommand{\poznamka}[1]{\par\begingroup\small\color{notegray}\itshape
  \textnormal{\textbf{Poznámka pro učitele.}}\enspace #1\par\endgroup}
\newcommand{\reseni}[1]{\par\begingroup\small\color{notegray}
  \textbf{Očekávané řešení:}\enspace #1\par\endgroup}
\newcommand{\blok}[1]{\needspace{5\baselineskip}\subsection*{#1}}
\newcommand{\vyrok}[1]{\uv{#1}}

\usepackage{tikz}
\nastavPripravu{1. ročník}{Soustavy rovnic}{79}{Geometrická interpretace soustav}
\begin{document}
\blok{Soustava jako průnik dvou přímek}
Každá rovnice $ax+by=c$ s alespoň jedním nenulovým koeficientem určuje v rovině přímku. Dvojice $(x;y)$ řeší soustavu právě tehdy, když leží na obou přímkách. Z hodin 77–78 známe výpočet; nyní algebraický výsledek spojíme s obrazem a počtem společných bodů.
\begin{definiceBox}
\textbf{Grafické řešení soustavy} je hledání společných bodů grafů obou rovnic v souřadnicové rovině. Množina řešení je průnik množin bodů obou přímek.
\end{definiceBox}
\begin{definiceBox}
Dvě přímky v rovině mohou mít \textbf{právě jeden společný bod}, být \textbf{rovnoběžné různé} (žádný společný bod), nebo \textbf{splývat} (nekonečně mnoho společných bodů). Pro soustavu tomu odpovídá jedno, žádné nebo nekonečně mnoho řešení.
\end{definiceBox}
\begin{prikladReseniBox}
Řešte početně a graficky soustavu $x+y=4$, $2x-y=5$.
\tcblower
Sečtením rovnic dostaneme $3x=9$, tedy $x=3$, a z $x+y=4$ vyjde $y=1$. Zkouška: $3+1=4$ a $2\cdot3-1=5$. Pro grafy vyjádříme $y=4-x$ a $y=2x-5$; první prochází $(0;4),(4;0)$, druhá např. $(2;-1),(3;1)$. Protínají se v $(3;1)$, tedy $K=\{(3;1)\}$.
\end{prikladReseniBox}
\begin{center}
\begin{tikzpicture}[x=.72cm,y=.63cm,>=stealth]
\draw[step=1,gray!25,very thin] (-1,-2) grid (5,5);
\draw[->] (-1.3,0)--(5.4,0) node[right] {$x$};\draw[->] (0,-2.3)--(0,5.4) node[above] {$y$};
\foreach \x in {1,2,3,4,5}{\draw (\x,.08)--(\x,-.08) node[below] {\small $\x$};}
\foreach \y in {1,2,3,4}{\draw (.08,\y)--(-.08,\y) node[left] {\small $\y$};}
\draw[very thick] (-1,5)--(5,-1);
\draw[very thick,dashed] (1.5,-2)--(5,5);
\fill (3,1) circle (2.7pt) node[above left=5pt] {\small $(3;1)$};
\end{tikzpicture}
\end{center}
\poznamka{Plná čára značí $y=4-x$, čárkovaná $y=2x-5$; jsou rozlišitelné i černobíle. Nech žáky nejprve odhadnout průsečík a pak vyžaduj přesnou zkoušku; z nepřesného náčrtu nemusí být možné odečíst zlomek.}
\blok{Dva krajní případy}
\begin{prikladReseniBox}
Porovnejte soustavy (A) $x+y=4$, $2x+2y=6$ a (B) $x+y=4$, $2x+2y=8$. Co ukáže graf a eliminace?
\tcblower
(A) Po odečtení dvojnásobku první rovnice od druhé vyjde $0=-2$, takže $K=\varnothing$. Grafy $y=4-x$ a $y=3-x$ mají stejnou směrnici $-1$, ale jiný úsek na ose $y$: jsou rovnoběžné různé. (B) Stejná eliminace dává $0=0$; druhá rovnice jen opakuje první. Grafy splývají a $K=\{(t;4-t):t\in\mathbb R\}$.
\end{prikladReseniBox}
\begin{center}
\begin{tikzpicture}[x=.58cm,y=.5cm,>=stealth]
\begin{scope}
\draw[step=1,gray!25,very thin] (-1,-1) grid (4,5);
\draw[->] (-1.2,0)--(4.3,0) node[right] {$x$};\draw[->] (0,-1.2)--(0,5.2) node[above] {$y$};
\draw[very thick] (-1,5)--(4,0);\draw[very thick,dashed] (-1,4)--(4,-1);
\node at (1.5,-2) {\small (A) žádné řešení};
\end{scope}
\begin{scope}[xshift=5.2cm]
\draw[step=1,gray!25,very thin] (-1,-1) grid (4,5);
\draw[->] (-1.2,0)--(4.3,0) node[right] {$x$};\draw[->] (0,-1.2)--(0,5.2) node[above] {$y$};
\draw[very thick] (-1,5)--(4,0);\draw[very thick,dashed] (-1,5)--(4,0);
\node at (1.5,-2) {\small (B) nekonečně mnoho};
\end{scope}
\end{tikzpicture}
\end{center}
\blok{Samostatné procvičení}
\begin{zadaniReseniBox}
\begin{enumerate}[label=\alph*)]
\item Sestrojte přímky $y=x+1$ a $y=-x+3$. Odečtěte jejich průsečík a ověřte jej výpočtem.
\item Určete bez kreslení počet řešení soustav $x-2y=3$, $2x-4y=6$ a $x-2y=3$, $2x-4y=7$. Vysvětlete polohu grafů.
\item Vymyslete druhou rovnici k $y=2x-1$ tak, aby soustava měla právě jedno řešení v $(1;1)$.
\end{enumerate}
\tcblower
\begin{enumerate}[label=\alph*)]
\item Rovnost $x+1=-x+3$ dává $2x=2$, $x=1$, $y=2$. Obě přímky procházejí $(1;2)$; první určují např. $(0;1),(1;2)$, druhou $(0;3),(1;2)$.
\item První druhá rovnice je dvojnásobkem první, grafy splývají a $K=\{(3+2t;t):t\in\mathbb R\}$. U druhé soustavy by dvojnásobek první pravé strany byl $6$, ale stojí $7$; přímky jsou rovnoběžné různé a $K=\varnothing$.
\item Např. $y=-x+2$ prochází $(1;1)$ a má jinou směrnici než $y=2x-1$, takže jde o jediný průsečík. Kontrola: $2\cdot1-1=1$ a $-1+2=1$.
\end{enumerate}
\end{zadaniReseniBox}
\newpage
\blok{Volitelné rozšíření a rezerva}
\begin{zadaniReseniBox}
Mohou se dvě různé přímky v rovině protnout ve dvou bodech? Jaký důsledek má odpověď pro počet řešení soustavy dvou nezávislých lineárních rovnic?
\tcblower
Ne. Dva různé body určují jedinou přímku; kdyby obě přímky obsahovaly stejné dva body, musely by být totožné. Dvě různé nerovnoběžné přímky se tedy protínají právě jednou a soustava má právě jednu dvojici.
\end{zadaniReseniBox}
\blok{Domácí práce}
\begin{zadaniReseniBox}
\begin{enumerate}[label=\alph*)]
\item Graficky i početně vyřešte $y=2x+1$, $y=-x+4$.
\item Určete počet řešení soustavy $3x+y=5$, $6x+2y=10$ a zapište celou množinu.
\item Změňte jen pravou stranu druhé rovnice v b) tak, aby nová soustava neměla řešení; zdůvodněte.
\end{enumerate}
\tcblower
\begin{enumerate}[label=\alph*)]
\item $2x+1=-x+4$ dává $x=1$, $y=3$. Grafy se protínají v $(1;3)$; zkouška $3=3$.
\item Druhá rovnice je dvojnásobkem první, tedy grafy splývají a $K=\{(t;5-3t):t\in\mathbb R\}$.
\item Např. $6x+2y=11$: levá strana je stále dvojnásobek první rovnice, jejíž dvojnásobná pravá strana je $10$. Vyšlo by $10=11$, tedy žádné společné řešení.
\end{enumerate}
\end{zadaniReseniBox}
\poznamka{Jádro: vazba soustavy na graf 8 min, průsečík a zkouška 15 min, tři polohy přímek 9 min, samostatné úlohy s kontrolou 10 min, organizace 3 min. Rezerva nad rámec; domácí práce asi 15 min. Navazující rovnice s absolutní hodnotou patří až hodině 80.}
\end{document}
